\paragraph{Charged phase-space data and neutral Gaussian packets.}
The charged trajectory can supply packet parameters only after its
positions and velocities are expressed in the Coulomb quotient. This is a
preparation construction in the full interacting Hilbert space of
Theorem~\ref{thm:whitney-interacting-quantum}. A family of prepared packets
at classical sample times does not constitute quantum propagation.

Let \(B\) have orthonormal columns spanning the mean-zero gauge space,
let \(T_c\) span \(\ker(D^{\mathsf T}M)\) orthonormally, and put
\(L_0=B^{\mathsf T}D^{\mathsf T}MDB\). For a differentiable
full configuration \(y(t)=(a(t),\psi(t))\), define
\begin{align}
 \xi&=-BL_0^{-1}B^{\mathsf T}D^{\mathsf T}Ma,&
 \dot\xi&=-BL_0^{-1}B^{\mathsf T}D^{\mathsf T}M\dot a,\nonumber\\
 a_C&=a+D\xi,& \psi_C&=e^{ie\xi}\psi,\nonumber\\
 \dot a_C&=\dot a+D\dot\xi,&
 \dot\psi_C&=e^{ie\xi}(\dot\psi+ie\dot\xi\mathbin{\odot}\psi).
 \label{eq:whitney-packet-coulomb-tangent}
\end{align}
Products and exponentials on nodal coefficients are componentwise.
The quotient coordinates are
\(q=(T_c^{\mathsf T}a_C,\operatorname{Re}\psi_C,
\operatorname{Im}\psi_C)\in\mathbb R^{56}\), with velocity \(w=\dot q\). The unit phase generator is
\(J(a,u,v)=(0,-v,u)\); its parameter is the phase angle,
without the charge factor \(e\).

\begin{proposition}[Coulomb phase-space preparation]
\label{prop:whitney-packet-coulomb}
For a temporal-gauge trajectory satisfying all thirteen Gauss equations,
the transformed scalar potential is \(\phi_C=-\dot\xi\), its constant
component vanishes, and the Schur minimizer is exactly \(-\dot\xi\).
The reduced cotangent is \(p=\gamma_qw\) and satisfies \(p(Jq)=0\).
For the recorded symmetric sector, with radial coefficient \(\alpha\),
\begin{equation}
 \xi=\frac{\alpha}{13}(12,-1,\ldots,-1),\qquad a_C=0.
 \label{eq:whitney-packet-radial-chart}
\end{equation}
At its exact initial datum, with \(\mathcal V=|\Omega|=10+10\sqrt5/3\),
\begin{equation}
 q_a=p_a=0,\quad \psi_C=\mathbf1,\quad
 p_{\operatorname{Re}\psi}=0,\quad
 p_{\operatorname{Im}\psi}
   =\mathcal V(3/10,-1/40,\ldots,-1/40).
 \label{eq:whitney-packet-initial-cotangent}
\end{equation}
\end{proposition}
\begin{proof}
The gauge law for the temporal covariant derivative is
\(\phi\mapsto\phi-\dot\xi\); differentiating the nodal gauge
law gives \eqref{eq:whitney-packet-coulomb-tangent}. The transformed
velocity plus vertical scalar-potential term is the image of the original
velocity under the time-independent gauge differential. Gauge covariance
therefore preserves every scalar-potential equation. Positivity of its
mean-zero Hessian makes \(-\dot\xi\) the unique minimizer. Since
\(\sum\xi_i=0\), the constant potential remains zero. The reduced
Legendre formula and constant Gauss equation give the cotangent and its
zero moment map. Without Gauss, the displayed rechart remains valid but
the equality of its transformed potential and Schur minimizer need not hold.

The radial incidence equation is \(\xi_j-\xi_0=-\alpha\),
and its mean-zero solution gives \eqref{eq:whitney-packet-radial-chart}.
Initially the edge dressing derivative vanishes because all nodal
coefficients equal one. Write \(A_0\) for the scalar nodal mass matrix.
The transformed covariant velocity is the original
\(i(3,-1,\ldots,-1)\), so the scalar cotangent is
\(2A_0i(3,-1,\ldots,-1)\). Simplex integration gives
\((A_0)_{00}=\mathcal V/10\),
\((A_0)_{0j}=\mathcal V/80\), and the boundary row sum
\(\mathcal V/16\). These identities prove
\eqref{eq:whitney-packet-initial-cotangent}. Maxwell orthogonality
\(T_c^{\mathsf T}MD=0\) gives the zero edge cotangent. Its constant
charge is \(\mathcal V(3/10-12/40)=0\).
\end{proof}

Supply a width \(\sigma>0\), momentum \(p\), center \(q\), and
\(\hbar>0\). In Lebesgue measure on all \(56\) quotient coordinates,
let the normalized seed be
\begin{equation}
 g_{q,p,\sigma}(x)=(2\pi\sigma^2)^{-14}
 \exp\!\left(-\frac{|x-q|^2}{4\sigma^2}
                  +\frac{i}{\hbar}p\cdot(x-q)\right).
 \label{eq:whitney-packet-seed}
\end{equation}
Its position covariance is \(\sigma^2I_{56}\); no transverse direction
has been removed. The unitary half-density map is
\(f\mapsto\sqrt\rho f\), where \(\rho=\sqrt{\det\gamma}\).
The circle acts orthogonally on the scalar coordinates and fixes the thirty
edge coordinates, so its average \(P_0\) has the same pullback expression
in both measures.

\begin{proposition}[A nonvanishing neutral projection]
\label{prop:whitney-packet-neutral-projection}
Write \(q_s,p_s\in\mathbb R^{26}\) for the scalar parts and set
\begin{equation}
 A=\frac{|q_s|^2}{4\sigma^2}+\frac{\sigma^2|p_s|^2}{\hbar^2},
 \qquad B_*=\frac{p(Jq)}{\hbar}.
 \label{eq:whitney-packet-overlap-parameters}
\end{equation}
Then \(A\geq|B_*|\) and
\begin{equation}
 n^2:=\|P_0g_{q,p,\sigma}\|^2
 =e^{-A}I_0\!\left(\sqrt{A^2-B_*^2}\right)>0.
 \label{eq:whitney-packet-neutral-norm}
\end{equation}
Consequently \(F=\rho^{-1/2}P_0g/n\) is a normalized vector in
\(\mathcal D(\widehat H_0)\). At the exact initial datum and
\(\hbar=1\), each width \(\sigma\in\{1/4,1/2,1\}\) satisfies
\begin{equation}
 B_*=0,\quad
 A=\frac{13}{4\sigma^2}+\frac{39\sigma^2\mathcal V^2}{400}<64,
 \qquad n^2>\frac1{64}.
 \label{eq:whitney-packet-initial-projection-bound}
\end{equation}
\end{proposition}
\begin{proof}
The Cauchy and arithmetic--geometric mean inequalities give
\(|B_*|\leq|p_s||q_s|/\hbar\leq A\). Completing the Gaussian square
in \(\langle g,\mathcal U_\theta g\rangle\) gives
\[
 \exp\bigl[-A(1-\cos\theta)-iB_*\sin\theta\bigr].
\]
Its circle average is \(n^2\). Expanding
\(\exp[(A-B_*)e^{i\theta}/2+(A+B_*)e^{-i\theta}/2]\)
and retaining equal powers gives
\(e^{-A}\sum_{k\geq0}((A^2-B_*^2)/4)^k/(k!)^2\), proving
\eqref{eq:whitney-packet-neutral-norm}. This also defines \(I_0\)
and proves positivity, including the equality case \(A=|B_*|\).
For \(B_*=0\) it is the usual real circle-integral representation.

Polynomial coefficient and density-derivative bounds from
Proposition~\ref{prop:whitney-interacting-gaussian-state} apply to every
translated, modulated Gaussian. Thus \(\rho^{-1/2}g\) belongs to
\(\mathcal D(\widehat H)\). The bounded circle projection preserves
that domain and commutes with the Hamiltonian. Division by \(n>0\)
proves the domain claim. The same argument gives finite moments of every
polynomial multiplication observable.

Initially \(|q_s|^2=13\) and
\(|p_s|^2=39\mathcal V^2/400\). Since \(\mathcal V<18\), direct
rational substitution at the three stated widths gives \(A<64\).
On \(|\theta|\leq1/8\),
\(A(1-\cos\theta)\leq A\theta^2/2<1/2\). Hence the integral is
larger than \(e^{-1/2}/(8\pi)>1/64\), using
\(e^{-1/2}>1/2\) and \(\pi<4\).
\end{proof}

Projection changes the state: a nonzero center is a parameter for an orbit,
not the scalar one-point expectation. In the Lebesgue half-density representation, circle invariance gives
\(\langle x_s\rangle=\langle -i\hbar\nabla_s\rangle=0\).
For \(B_*=0\), direct Gaussian integration gives the useful invariant
nodal-radius observable
\begin{equation}
 \langle|x_s|^2\rangle_F
 =26\sigma^2+\frac{|q_s|^2}{2}(1+r)
       -\frac{2\sigma^4|p_s|^2}{\hbar^2}(1-r),
 \qquad r=\frac{I_1(A)}{I_0(A)}.
 \label{eq:whitney-packet-radius-moment}
\end{equation}
Indeed, inserting \(|x_s|^2\) in the overlap adds
\(26\sigma^2+|q_s|^2(1+\cos\theta)/2
-2\sigma^4|p_s|^2(1-\cos\theta)/\hbar^2\); the imaginary term
is proportional to \(B_*\). The cosine average is \(I_1/I_0\). For arbitrary \(B_*\), the same
calculation gives
\(26\sigma^2+|q_s|^2/2-2\sigma^4|p_s|^2/\hbar^2
+2\sigma^2zI_1(z)/I_0(z)\), where
\(z=\sqrt{A^2-B_*^2}\); its last term is zero when \(z=0\).
The thirty edge coordinates retain their seed covariance, independently
of this scalar projection. This nodal observable is not a supplied
laboratory measurement operator.

\paragraph{The residual required for propagation.}
The half-density change of measure does not flatten the Hamiltonian.
Writing \(a^{ij}=(\gamma^{-1})^{ij}\) and \(\ell=\log\rho\), its
Lebesgue expression is
\begin{equation}
 \widetilde H=-\frac{\hbar^2}{2}\partial_i(a^{ij}\partial_j)+V+U_\rho,
 \qquad U_\rho=\frac{\hbar^2}{4}\partial_i(a^{ij}\partial_j\ell)
      +\frac{\hbar^2}{8}(\partial_i\ell)a^{ij}(\partial_j\ell).
 \label{eq:whitney-packet-half-density-operator}
\end{equation}
This follows by differentiating \(\rho^{-1/2}g\); the two first-order
cross terms cancel by symmetry of \(a^{ij}\). For any smooth seed path
in the strong domain, let \(r_g=i\hbar\dot g-\widetilde Hg\).
Where \(n(t)=\|P_0g(t)\|>0\), the normalized projection has residual
norm at most \(2\|r_g(t)\|/n(t)\): differentiating the projected norm
and using self-adjointness gives
\(|\dot n|\leq\|P_0r_g\|/\hbar\). Duhamel's formula therefore gives
\begin{equation}
 \left\|e^{-it\widehat H_0/\hbar}F(0)-F(t)\right\|
 \leq\frac2\hbar\int_0^{|t|}\frac{\|r_g(s)\|}{n(s)}\,ds
 \label{eq:whitney-packet-projected-residual}
\end{equation}
for forward time, and with the reversed path for negative time.
The initial bound on \(n\) alone does not bound this integral. Continuous
center and covariance defects, the full variable coefficients and
\(U_\rho\), their numerical errors, and all Gaussian tails must enter
an actual residual certificate. The prepared samples supply none of those
propagation estimates. In particular, no error bound on a \(1/40\)
interval at \(\hbar=1\) follows from the preparation or from a
semiclassical limit.

The executable packet at
\path{code/electromagnetism/runtime/whitney_quantum_packet_receipt.json}
recharts two adjacent classical samples and checks all mean-zero multipliers
and reduced cotangents. Its independent checker uses exact-degree simplex
moments at \(a_C=0\), including the complete \(68\)-coordinate kinetic
metric and all \(56\) reduced directions. Matrix operations and the second
classical sample are numerical, while the displayed initial cotangent and
projection lower bound are analytic identities. The packet specifies three
full-dimensional seeds and their neutral projections. It does not select a
physical preparation, evolve their covariances, identify them with a computed
quantum history, or establish interacting continuum field theory.
